\section{Limits and What We Did Not Test}
\label{sec:limits}

Our simulator runs as a single program, so all shards share one machine.
It does not model network delays, waiting in queues, work done at the same time, or failures.
All runs use one data set and one text model, with \RNumCorpusTotalPassages{} passages drawn once with a fixed seed.
They also use five seeds and fixed sizes: \RNumNumShards{} shards, two searched shards per query, and \RNumCentersPerSummary{} centers per summary.

About \RNumSteadyEpochInsertVecs{} passages arrive per time step, and a summary stays the same size as arrivals grow.
But the share of new passages that push can keep at a fixed budget falls, so higher arrival rates may favor refresh.
In one-shard and eight-shard placement, new passages join shards that keep their starting passages, and a fresh summary covers the whole shard.
On mixed topics this lowered overall recall, plausibly because the centers shift toward the random new passages.
Separate summaries for new data might avoid this, but we did not model them.
A coordinator that knows which shards are open could also reach them without upkeep traffic.
Saving re-sent summary bytes, carrying unused bytes over, and a refresh rule based on the plain count are also untested.
